\section{Methodology}

% \subsection{Dataset Preparation}
% Although the public IPMet files are described in the dataset documentation as grayscale radar reflectivity images, direct inspection showed that the downloaded PNG files are stored in RGBA format. The alpha channel contains both fully and partially transparent pixels; therefore, simply discarding this channel would not provide a consistent representation of the transparent background regions. To obtain a uniform three-channel representation suitable for the heterogeneous super-resolution models, each RGBA image was alpha-composited onto a fixed neutral-gray background with an intensity value of 128. The resulting RGB images were retained as the input representation for the super-resolution experiments, thereby preserving the available color and structural information of the radar imagery. The original alpha channel was stored separately for subsequent characterization of echo coverage and spatial structure. Since the public PNG files do not provide an explicit mapping between the rendered RGB values and physical dBZ reflectivity, image-derived intensity characteristics were not interpreted directly as physical reflectivity values.

% Temporal Data Partitioning. Radar observations acquired within short temporal intervals may contain highly similar precipitation structures, and conventional random image-level partitioning can therefore introduce information leakage between training and evaluation sets. To reduce this risk, the dataset was partitioned using temporal acquisition blocks rather than randomly selected individual images. After removal of the exact duplicate identified during dataset inspection, the remaining 4,013 images were chronologically ordered according to their acquisition timestamps. A new temporal block was initiated whenever the interval between consecutive radar observations exceeded 60 minutes. This procedure produced 31 non-overlapping temporal acquisition blocks. Entire blocks were then assigned chronologically to the training, validation, and test partitions while approximating a 70/15/15 distribution. The resulting partitions contained 2,799 training images from 18 blocks (69.75%), 604 validation images from 7 blocks (15.05%), and 610 test images from 6 blocks (15.20%). No temporal block was shared between partitions. The temporal separation between the final training observation and the first validation observation was 128 minutes, while the separation between validation and test was 713 minutes. The validation partition is reserved for model selection, radar-regime characterization, coverage-threshold estimation, and ensemble pruning, whereas the test partition is retained exclusively for final evaluation.

% Super-Resolution Degradation. Following preprocessing, the resulting RGB radar images were used as high-resolution reference images. Low-resolution counterparts were generated through bicubic downsampling using scale factors of ×2, ×3, and ×4, representing progressively stronger degradation conditions. These scales were selected to provide multiple super-resolution difficulty levels and are consistent with degradation factors previously evaluated on IPMet radar imagery. To ensure exact spatial correspondence across all three scaling factors, the learning pipeline operates on 192 × 192 high-resolution patches, producing native low-resolution patches of 96 × 96, 64 × 64, and 48 × 48 pixels for ×2, ×3, and ×4, respectively. The original high-resolution RGB patch serves as the reconstruction target. Native low-resolution inputs are retained for architectures that perform internal upsampling, while bicubically upsampled counterparts can be supplied to architectures designed for same-resolution restoration. This common degradation protocol is applied consistently across all members of the heterogeneous model pool.


% Patch Construction. Following temporal partitioning, the 4,013 radar images were spatially decomposed into overlapping high-resolution patches of 192 × 192 pixels. This patch size was selected because it is exactly divisible by the three considered super-resolution scale factors (×2, ×3, and ×4), resulting in corresponding low-resolution dimensions of 96 × 96, 64 × 64, and 48 × 48 pixels, respectively. Patches were extracted using a nominal stride of 128 pixels, with an additional boundary-aligned position included whenever the regular stride did not reach the image boundary. For the 640 × 640 IPMet images, this produced five horizontal and five vertical positions and therefore 25 patches per image. In total, 100,325 patches were indexed, comprising 69,975 training, 15,100 validation, and 15,250 test patches. Rather than storing cropped images separately, patch coordinates and source-image metadata were retained in a manifest to enable dynamic extraction during model training and evaluation. No patches were discarded according to precipitation sparsity, since sparse conditions constitute an important radar regime in the proposed ensemble-selection framework. Auxiliary patch descriptors, including alpha-visible content fraction and normalized distance from the radar centre, were also retained for subsequent regime characterization.

% Data Loading and Multi-Scale Degradation. A PyTorch-based data pipeline was developed to dynamically generate super-resolution input–target pairs from the indexed radar patches. Each high-resolution patch was represented as a three-channel RGB tensor of size 192 × 192 and normalized to the [0,1] range. For each super-resolution scale, bicubic downsampling was applied to generate native low-resolution patches of 96 × 96, 64 × 64, and 48 × 48 pixels for ×2, ×3, and ×4 degradation, respectively. The corresponding high-resolution 192 × 192 patch was retained as the reconstruction target. In addition, each native low-resolution patch was bicubically resized back to 192 × 192 to support architectures that operate on same-resolution degraded inputs. The original alpha patch and patch-level metadata were preserved alongside each sample for later radar-regime characterization and meteorology-aware evaluation. Separate fixed-scale validation and test datasets were maintained for ×2, ×3, and ×4 to enable degradation-specific performance analysis.


% Data Loading and Batch Construction. Separate PyTorch DataLoaders were constructed for the ×2, ×3, and ×4 degradation levels to accommodate the different native low-resolution spatial dimensions associated with each scale. For a batch size of 16, the resulting low-resolution tensors had dimensions of 96 × 96, 64 × 64, and 48 × 48 pixels for ×2, ×3, and ×4, respectively, while the corresponding high-resolution targets remained fixed at 192 × 192 pixels. The training partition produced 4,374 batches per scale, whereas the validation and test partitions produced 944 and 954 batches per scale, respectively. Each batch also retained the associated alpha mask and patch metadata to support subsequent radar-regime characterization and meteorology-aware evaluation. Separate scale-specific loaders were maintained throughout training and evaluation so that performance differences associated with degradation severity could be explicitly analyzed.